\documentclass[11pt]{article}
\usepackage[margin=1in]{geometry}
\usepackage{booktabs}
\usepackage{array}
\usepackage{hyperref}
\usepackage{listings}

\title{Verification-Guided RTL Repair Agent for Out-of-Order Processor Bugs}
\author{Vayun Gupta}
\date{}

\begin{document}
\maketitle

\section*{Question}
This project studies a practical hardware-debug question: when an LLM is asked
to repair broken RTL from a real processor codebase, which debug artifacts make
the repair more reliable? The controlled comparison is between three evidence
conditions: source-level RTL context only, RTL plus simulator failure logs, and
RTL plus simulator logs plus a short waveform-style signal summary.

This matters because real RTL debugging is rarely a clean code-generation task.
An engineer usually starts from a failing simulation, a failing assertion, a
waveform, and a large set of plausible files. The useful question is therefore
not whether an LLM can write a module from scratch, but whether realistic
verification evidence helps it localize and repair the kind of temporal bugs
that appear in processor development.

\section*{Prior Prediction}
Before running the substantial experiments, I wrote the dated prediction file
\texttt{predictions.md} on 2026-05-14. The primary prediction was that simulator
logs would improve repair success over source-only repair, and that adding
waveform-derived signal summaries would help most on temporal or protocol-level
bugs. In particular, I expected the trace summaries to matter most when the root
cause was a relation across cycles, such as a same-cycle rename bypass, a
multi-beat memory transaction, or a branch-recovery state transition.

\section*{System Overview}
The project implements a small verification-guided RTL repair benchmark and
agent loop. The pipeline is:

\begin{center}
\small
\begin{tabular}{p{0.20\linewidth}p{0.68\linewidth}}
\toprule
Stage & Engineering role \\
\midrule
Bug mining & Inspect ECE 411 processor commit history and identify real bug-fix commits with local, testable RTL behavior. \\
Bug injection & Start from the current known-good Verilator-compatible RTL and apply a targeted mutation that recreates the historical bug. \\
Validation & Run the reference and buggy versions through the HW4-style verifier; keep a case only if reference passes and buggy fails. \\
Artifact generation & Produce simulator failure logs and a concise trace summary that names the relevant signal relationship. \\
Repair & Ask the model to patch the RTL under one evidence condition. In localization mode, it sees multiple plausible files. \\
Verification & Compile and run the repaired module against the hidden Verilator testbench and record pass/fail JSON logs. \\
\bottomrule
\end{tabular}
\end{center}

The implementation reuses the HW4 verifier interface and per-trial JSON format,
but the problem suite is new. The final-project code is organized as follows:
\texttt{bug\_cases.py} defines the injected bugs and testbenches;
\texttt{artifacts.py} builds the debug artifacts; \texttt{repair\_agent.py}
constructs the prompts and checks model patches; and \texttt{evaluate.py} runs
the controlled study and writes auditable per-trial logs.

\section*{Benchmark Construction}
The most industry-relevant part of the project is the benchmark construction.
Rather than relying only on synthetic mistakes, I mined a private ECE 411
out-of-order processor repository for bug-fix commits, then converted selected
fixes into controlled repair cases. A case was included only when the bug could
be isolated to a small RTL behavior and tested under Verilator without relying
on the original Synopsys environment.

\begin{center}
\small
\begin{tabular}{p{0.08\linewidth}p{0.15\linewidth}p{0.57\linewidth}}
\toprule
Case & Commit & Recreated failure \\
\midrule
C1 & \texttt{bbb61a8} & Instruction burst adapter starts a memory transaction even when the instruction cache lost arbitration. \\
C2 & \texttt{513f159} & RAT read path misses a same-cycle speculative destination bypass. \\
C3 & \texttt{582bf6c} & Dispatch observes a stale registered ROB index instead of the current tail index. \\
C4 & \texttt{2c76489} & Branch recovery restores a speculative freelist checkpoint instead of the committed shadow state. \\
C5 & \texttt{5ae5544} & JALR target update masks the immediate but omits the resolved \texttt{rs1} base. \\
\bottomrule
\end{tabular}
\end{center}

These real bugs are supplemented by seven synthetic stress cases that cover
similar RTL-debug patterns: FIFO flush priority, payload integrity, multi-beat
burst timing and packing, RAT commit/flush interaction, freelist resource
accounting, and ROB completion ordering. The main result in this report focuses
on the five real-history cases because they are the best match to the final
project goal of moving beyond a clean benchmark toward realistic engineering
use.

\section*{Controlled Comparison}
The main experiment uses the real-history cases in fault-localization mode. For
each bug, each condition is run for three trials, giving 45 total repair
attempts. The only intended difference between conditions is the evidence given
to the model:

\begin{center}
\small
\begin{tabular}{p{0.18\linewidth}p{0.68\linewidth}}
\toprule
Condition & Evidence available to the repair agent \\
\midrule
\texttt{rtl\_only} & Candidate module interfaces and the need to repair RTL, but no failure reason, design intent, or trace clue. \\
\texttt{sim\_log} & Problem statement, success criteria, candidate RTL files, and the simulator failure log. \\
\texttt{sim\_log\_trace} & Everything in \texttt{sim\_log}, plus a short waveform-style summary of the relevant signal timeline. \\
\bottomrule
\end{tabular}
\end{center}

This is not a claim that the three conditions are equally informative to a human
engineer. They are intentionally different information budgets. The goal is to
measure whether adding realistic verification artifacts changes the model's
ability to repair temporal RTL bugs.

\section*{Validation}
The verifier is the main guardrail against confusing execution with success.
Every case is checked in two directions before being used: the reference RTL
must pass the hidden testbench, and the injected buggy RTL must fail with the
expected symptom. The current real-case validation artifact,
\texttt{real\_case\_validation.json}, shows 5/5 references passing and 5/5
buggy versions failing.

The repair metric is strict: a trial counts as a pass only if the model-produced
SystemVerilog compiles and passes the hidden Verilator testbench. Syntax errors,
interface mismatches, partial patches, and behaviorally wrong repairs all count
as failures. The per-trial results are stored in \texttt{repair\_results.json},
including the case id, condition, trial number, failure reason, model
confidence, and prompt size.

\section*{Results}
The trace-summary condition was the strongest. Across the five real-history
bugs and three trials per condition, the source-only baseline repaired 6 of 15
trials, simulator logs repaired 8 of 15, and simulator logs plus trace summaries
repaired all 15.

\begin{center}
\begin{tabular}{lcc}
\toprule
Condition & Passes / Total & Pass rate \\
\midrule
RTL only & 6 / 15 & 40.0\% \\
Simulation log & 8 / 15 & 53.3\% \\
Simulation log + trace summary & 15 / 15 & 100.0\% \\
\bottomrule
\end{tabular}
\end{center}

\begin{center}
\small
\begin{tabular}{lccc}
\toprule
Case & RTL only & Sim log & Log + trace \\
\midrule
C1: i-cache arbitration & 3/3 & 3/3 & 3/3 \\
C2: RAT rename bypass & 0/3 & 0/3 & 3/3 \\
C3: ROB dispatch index & 3/3 & 3/3 & 3/3 \\
C4: freelist recovery & 0/3 & 0/3 & 3/3 \\
C5: JALR target & 0/3 & 2/3 & 3/3 \\
\bottomrule
\end{tabular}
\end{center}

\section*{Discussion}
The result supports the prior prediction, but with an important nuance. Two
bugs were easy in every condition: the instruction-cache arbitration bug and
the stale ROB dispatch index. These failures are local enough that the model can
often infer the patch from the interface and nearby RTL structure alone.

The more interesting cases were the temporal state bugs. The RAT bypass and
freelist recovery bugs both failed all source-only and log-only trials, but
passed all trace-summary trials. In these cases, the simulator failure reason
named the symptom, but not the state relationship that had to be preserved. The
trace summary made the missing invariant explicit: same-cycle speculative
rename mapping must bypass the RAT table, and branch recovery must restore the
freelist from committed shadow state rather than a stale speculative
checkpoint.

The JALR target bug sits between those extremes. Simulator logs were enough in
two of three trials because the failure message pointed at the target address,
but trace summaries made the repair consistent across all trials by spelling out
the relation between the stored immediate, the resolved \texttt{rs1} base, and
the final masked target.

From an engineering perspective, this suggests that the valuable artifact is
not merely ``more text.'' The useful artifact is a compact statement of the
cycle-level invariant that the waveform demonstrates. That is close to how a
hardware engineer uses a waveform: not as a dump of every signal, but as
evidence for a small causal story about state, timing, and protocol.

\section*{Limitations and Threats to Validity}
The strongest argument against the conclusion is sample size. The main
real-history study has five bugs, which fits the guideline's recommended
5--20 carefully chosen cases, but it is still a curated subset of one processor
project. The result should therefore be read as a controlled characterization,
not a universal claim that trace summaries always make RTL repair succeed.

The trace summaries are hand-authored from targeted testbench behavior rather
than automatically extracted from VCD files. This is useful for a clean course
experiment because the evidence is auditable, but it understates the difficulty
of building a production waveform summarizer. A deployable version would need
to extract candidate signal timelines automatically and avoid leaking the repair
in overly direct language.

Another limitation is that the repair model sees a small set of candidate files
in localization mode, not an entire industrial RTL tree. The benchmark is more
realistic than exact-module reconstruction, but still much smaller than a
commercial verification environment. Finally, the pass/fail verifier checks the
targeted hidden behavior; it is possible for a patch to pass this test while
remaining incomplete for the full processor.

\section*{Reproducibility}
The project includes the source code, dated prediction file, generated result
JSON, and validation artifacts required to audit the comparison. No API keys are
stored in the submitted files; the OpenAI key is read from the environment.

To validate the bug suite:
\begin{lstlisting}[basicstyle=\ttfamily\small]
cd FinalProject
python3 evaluate.py --validate-cases-only --real-only \
  --output-json real_case_validation.json
\end{lstlisting}

To reproduce the main real-history fault-localization study:
\begin{lstlisting}[basicstyle=\ttfamily\small]
python3 evaluate.py --real-only --localize --resume --trials 3 \
  --conditions rtl_only sim_log sim_log_trace \
  --output-json repair_results.json
\end{lstlisting}

The expected output is a JSON file with 45 real-history trials and aggregate
condition results matching the table above. The main known failure mode is API
token-rate pressure on long prompts; \texttt{evaluate.py} supports
\texttt{--resume} and prompt shrinking so interrupted runs can continue without
rerunning completed trials.

\end{document}
